% Section 3.2, from lectures/L03/README.md: what you should be able to do, the questions to test
% yourself with, and the next lecture.

\section{Sammanfattning}
\label{c3:sec:summary}

\needspace{6\baselineskip}
\subsection*{Det här ska du kunna nu}
Efter det här kapitlet och Labb~1 ska du kunna:
\begin{itemize}
  \item Avgöra vilka anslutningar på en tryckknapp som är slutande och vilka som är brytande, både
    med hjälp av märkningen och genom mätning.
  \item Koppla upp ett kontaktnät efter ett strömvägsschema, och kontrollera det mot en
    sanningstabell.
  \item Realisera AND, OR, NOT, NAND, NOR, XOR och XNOR med kontakter.
  \item Ta fram ett kontaktnät ur en sanningstabell, via ett förenklat uttryck.
  \item Analysera ett kontaktnät: ta fram dess uttryck och sanningstabell, och hitta ett mindre nät
    med samma funktion.
  \item Felsöka ett kontaktnät systematiskt, med en lampa eller en multimeter.
\end{itemize}

\needspace{6\baselineskip}
\subsection*{Frågor att testa dig själv med}
\begin{enumerate}
  \item Varför ska ett kontaktnät kopplas med spänningen frånslagen, fast 24~V inte är farligt att
    röra?
  \item Hur avgör du med en lampa och ett aggregat om ett kontaktelement är slutande eller brytande?
  \item Två kontaktnät ser helt olika ut. Hur avgör du om de gör samma sak?
  \item I \hyperref[lab1:u8]{uppgift~8} i Labb~1 går det att ta bort en kontakt utan att funktionen
    ändras. Hur kan en kontakt vara onödig, när den ju faktiskt påverkar strömmen?
  \item Lampan lyser aldrig, vad du än trycker på. Var börjar du leta, och varför just där?
  \item Varför används en stoppknapps brytande kontakt, och inte dess slutande, i en maskins
    stoppkrets?
\end{enumerate}

\needspace{6\baselineskip}
\subsection*{Nästa kapitel}
Kapitel~\ref{c4:ch} handlar om Karnaughdiagram och IC-kretsar:
\begin{itemize}
  \item Karnaughdiagram: ett systematiskt sätt att hitta det minsta nätet, i stället för att leta
    med räknelagarna.
  \item Syntes av kombinatoriska nät, från specifikation till grindnät.
  \item NAND- och NOR-nät.
  \item IC-kretsar ur 74-serien: benplacering, datablad och kopplingsdäck, inför Labb~2.
\end{itemize}
